\chapter{Supercomputing Deployment and Systems Engineering}
\label{ch:supercomputing}

Training custom neural networks and running inference on over 13,000 scanned exam pages required significant computational horsepower. In this project, all model training, evaluation, and batch inference runs were carried out on the \textbf{PARAM Kamrupa Supercomputer}, an institutional tier-1 supercomputing facility located at the Computer \& Communication Centre (CCC) at IIT Guwahati.

Working on a shared production supercomputer is very different from developing on a local laptop. Compute nodes are completely air-gapped without internet access, the operating system uses legacy enterprise libraries, and the GPU hardware imposes strict memory limits. In this chapter, we outline the hardware setup on PARAM Kamrupa, explain how we provisioned model weights offline, describe our workaround for legacy CentOS 7 dynamic linking issues, and detail how we tuned attention kernels and image resolution to fit a 7.6-billion parameter model into a 16GB VRAM budget.

\section{Supercomputing Cluster Architecture (PARAM Kamrupa)}
\label{sec:hpc_architecture}

PARAM Kamrupa was deployed under the National Supercomputing Mission (NSM) by C-DAC and IIT Guwahati. The cluster is built on the BullSequana X1000 blade server platform. Table \ref{tab:param_kamrupa_specs} summarizes the exact hardware specifications of the GPU partition we used.

\begin{table}[htbp]
\centering
\caption{PARAM Kamrupa Supercomputer GPU Node Hardware Specifications}
\label{tab:param_kamrupa_specs}
\vspace{0.2cm}
\small
\begin{tabular}{ll}
\toprule
\textbf{System Parameter} & \textbf{Hardware Specification} \\
\midrule
Cluster Platform & BullSequana X1000 Blade Platform \\
GPU Accelerators & \textbf{NVIDIA Tesla V100-SXM2} (Volta Architecture) \\
GPU Memory Capacity & \textbf{16 GB HBM2} (High-Bandwidth Memory @ 900 GB/s) \\
Compute Capability & \textbf{sm\_70} (Volta Microarchitecture) \\
CUDA / Tensor Cores & 4,096 CUDA Cores / 640 Tensor Cores per GPU \\
Host Processor & Dual-Socket Intel Xeon Gold 6248 (40 Cores / 80 Threads) \\
Host RAM & 192 GB -- 384 GB DDR4 Registered ECC \\
Storage Architecture & Parallel \textbf{Lustre Filesystem} \\
Scratch Workspace & \texttt{/scratch/p.dinkesh/hmer\_project/} \\
Cluster Interconnect & Mellanox InfiniBand EDR (100 Gbps) \\
Operating System & Red Hat Enterprise Linux 7.9 / CentOS 7 \\
Workload Orchestrator & SLURM (Simple Linux Utility for Resource Management) \\
\bottomrule
\end{tabular}
\end{table}

The `gpu` partition consists of compute nodes (`gpu001` through `gpu010`). Each node has four Tesla V100 GPUs connected via NVLink (300 GB/s bidirectional bandwidth), all managed through the SLURM queue.

\section{Air-Gapped Execution and Model Weight Caching}
\label{sec:air_gapped_provisioning}

For institutional security, \textbf{the compute nodes on PARAM Kamrupa have no outbound internet connection}. They cannot make DNS requests or download packages from PyPI, GitHub, or Hugging Face.

\subsection{The Offline Loading Problem}
Modern deep learning libraries often assume they can connect to the internet to download tokenizer configs and model weights on the fly. Running an off-the-shelf Hugging Face script inside a SLURM batch job immediately throws a network timeout error. Furthermore, our personal home directories have a strict 20 GB quota, which is not enough to store multiple 15 GB model checkpoints.

\subsection{Our Lustre Scratch Provisioning Workflow}
To solve this, we wrote an offline caching script (\texttt{src/ocr/download\_vlm\_weights.py}):
\begin{enumerate}
    \item \textbf{Downloading via Login Nodes:} Login nodes (`login01`--`login04`) have proxy internet access. We ran our download scripts on a login node to pull the raw model weights, configuration files, and tokenizers.
    \item \textbf{Storing on Lustre Scratch:} We used \texttt{huggingface\_hub.snapshot\_download} to save model files directly into the high-speed Lustre parallel filesystem at \texttt{/scratch/p.dinkesh/hmer\_project/models/qwen\_vl}, skipping unnecessary file types (\texttt{*.msgpack}, \texttt{*.h5}).
    \item \textbf{Shared Python Virtual Environment:} We built a self-contained virtual environment inside \texttt{/scratch/p.dinkesh/hmer\_project/venv/}, pre-installing PyTorch 2.0.1, Torchvision 0.15.2, and Transformers 4.46.3.
\end{enumerate}
When our batch jobs run on compute nodes, they load the model locally from the Lustre scratch drive at over $2.5\,\text{GB/s}$ across InfiniBand, without needing any internet connection.

\section{Fixing CentOS 7 C++ Library Mismatches via \texttt{ctypes}}
\label{sec:cxxabi_resolution}

One of the trickiest systems bugs we encountered on PARAM Kamrupa was an operating system library incompatibility. The cluster runs CentOS 7, which still uses an old GNU C++ compiler (\texttt{GCC 4.8.5}).

\subsection{The CXXABI Version Mismatch}
The default C++ runtime on CentOS 7 (\texttt{/lib64/libstdc++.so.6}) only supports symbols up to \texttt{GLIBCXX\_3.4.19} and \texttt{CXXABI\_1.3.7}. But modern PyTorch and Torchvision binaries (specifically the spatial operators in \texttt{torchvision.ops}) require \texttt{CXXABI\_1.3.9} and \texttt{GLIBCXX\_3.4.21}, which were introduced in GCC 5.1+. When we tried to import \texttt{torchvision} on a compute node, Python crashed with:
\begin{verbatim}
ImportError: /lib64/libstdc++.so.6: version `CXXABI_1.3.9' not found 
(required by .../torchvision/image.so)
\end{verbatim}
Because we do not have root privileges on the supercomputer, we could not update system packages using \texttt{yum}.

\subsection{Runtime Preloading via \texttt{ctypes}}
Instead of asking cluster administrators to rebuild the OS libraries, we found a clean software workaround. We noticed that the Anaconda installation on the cluster (\texttt{/opt/ohpc/pub/apps/Anaconda3-2023.03/}) included a modern version of \texttt{libstdc++.so.6}. 

We added a preloading hook using Python's \texttt{ctypes} library at the very top of our inference and training entry points (\texttt{run\_full\_page\_inference.py} and \texttt{download\_vlm\_weights.py}):
\begin{equation}
\label{eq:ctypes_preload}
\begin{aligned}
\texttt{ctypes.CDLL}(&\texttt{"/opt/ohpc/pub/apps/Anaconda3-2023.03/lib/libstdc++.so.6"}, \\
&\texttt{mode} = \texttt{ctypes.RTLD\_GLOBAL})
\end{aligned}
\end{equation}
By loading the newer C++ library into the global symbol table (\texttt{RTLD\_GLOBAL}) before any machine learning packages are loaded, all subsequent requests for \texttt{CXXABI\_1.3.9} resolve against Anaconda's runtime rather than the outdated CentOS 7 library. This completely eliminated the import crashes across all compute nodes.

\section{Volta Architecture (sm\_70) Attention Kernels and VRAM Limits}
\label{sec:volta_vram_optimization}

The Tesla V100 GPU is based on the **Volta microarchitecture (sm\_70)**. Running a 7.6-billion parameter vision-language model on Volta hardware requires careful memory budgeting.

\subsection{Why FlashAttention-2 Fails on Volta and Cutlass SDPA Fallback}
Most modern LLM inference engines use **FlashAttention-2** \cite{dao2022flashattention} to save memory. However, FlashAttention-2 relies on special asynchronous copy instructions (\texttt{cp.async}) that are physically present only on newer NVIDIA architectures (Ampere sm\_80, Ada sm\_89, and Hopper sm\_90). When we tried to run FlashAttention-2 on our Tesla V100 (sm\_70), kernel compilation failed immediately.

We fixed this by configuring PyTorch's native Scaled Dot-Product Attention (SDPA) to use the **Cutlass Memory-Efficient Attention** kernel instead:
\begin{align}
\texttt{torch.backends.cuda.enable\_flash\_sdp(False)} \\
\texttt{torch.backends.cuda.enable\_mem\_efficient\_sdp(True)} \\
\texttt{torch.backends.cuda.enable\_math\_sdp(True)}
\end{align}
Cutlass computes exact attention without constructing the full $N \times N$ attention matrix in memory, keeping attention memory bounded to $\mathcal{O}(N)$ on Volta GPUs.

\subsection{Fitting a 7.6B Model into 16GB VRAM}
Our flagship model, **Qwen2.5-VL-7B-Instruct**, contains about $7.61 \times 10^9$ parameters. In 16-bit half-precision (\texttt{float16}), loading the weights alone takes **$14.07\,\text{GiB}$** of GPU memory. Since the Tesla V100-PCIe partition on PARAM Kamrupa has an addressable limit of **$15.78\,\text{GiB}$** ($16,160\,\text{MiB}$), the remaining headroom for CUDA context, activations, KV caches, and visual tokens is exceedingly tight:
\begin{equation}
M_{\text{headroom}} = 15.78\,\text{GiB} - 14.07\,\text{GiB} \approx \mathbf{1.71\,\text{GiB}} \quad (1,750\,\text{MiB}).
\end{equation}

Table \ref{tab:vram_arithmetic} details the exact VRAM allocation arithmetic that distinguishes monolithic full-page processing from our multi-band slicing approach.

\begin{table}[htbp]
\centering
\caption{VRAM Allocation Breakdown on PARAM Kamrupa NVIDIA Tesla V100 (16GB)}
\label{tab:vram_arithmetic}
\vspace{0.2cm}
\small
\begin{tabular}{p{4.6cm}p{2.6cm}p{2.4cm}p{3.6cm}}
\toprule
\textbf{Memory Allocation Component} & \textbf{Monolithic Full-Page} & \textbf{3-Band Slicing} & \textbf{Operational Impact} \\
\midrule
Total Addressable VRAM & $15.78\,\text{GiB}$ & $15.78\,\text{GiB}$ & Physical hardware ceiling \\
Qwen2.5-VL-7B FP16 Weights & $14.07\,\text{GiB}$ & $14.07\,\text{GiB}$ & Static model footprint \\
PyTorch CUDA Context / Buffers & $350\,\text{MiB}$ & $350\,\text{MiB}$ & Base runtime overhead \\
\midrule
\textbf{Available Headroom} & \textbf{1,750 MiB} & \textbf{1,750 MiB} & Addressable buffer \\
\midrule
ViT Input Image Patches ($N$) & $> 5,800$ patches & $\approx 1,900$ patches & $67\%$ sequence reduction \\
ViT Self-Attention Contiguous Tensor & \textbf{1.30 GiB (1,331 MiB)} & \textbf{240 MiB} & $81.5\%$ activation drop \\
Allocator Memory Fragmentation & High ($> 400\,\text{MiB}$) & Low ($< 80\,\text{MiB}$) & Contiguous block availability \\
\midrule
\textbf{Final Safety Margin} & \textbf{< 419 MiB} & \textbf{> 1,430 MiB} & Zero vs. constant crash risk \\
\textbf{Execution Result} & \textbf{CUDA OOM Crash} & \textbf{Stable Execution} & 100\% execution reliability \\
\bottomrule
\end{tabular}
\end{table}

\subsection{Dynamic Image Resolution and Multi-Band Scaling}
As established in Table \ref{tab:vram_arithmetic}, attempting to feed an uncompressed full-page scan ($2338 \times 1653$) directly forces the attention tensor to allocate $1.30\,\text{GiB}$ contiguously. Because memory fragmentation prevents PyTorch's caching allocator from finding a single contiguous $1.30\,\text{GiB}$ chunk in the $1.75\,\text{GiB}$ headroom, the process terminates immediately with an Out-of-Memory error.

While downscaling the entire page to $848\,\text{px}$ allowed the process to fit, it caused unacceptable visual blurring on subscripts and minus signs. In contrast, our \textbf{Multi-Band High-Resolution Slicing Architecture} resolves this dilemma cleanly: by feeding horizontal bands at $\text{max\_dim} = 1200\,\text{px}$, each band generates only $\approx 1,900$ patches. The ViT attention memory collapses to just $240\,\text{MiB}$, leaving over $1.43\,\text{GiB}$ of safety headroom while preserving $100\%$ optical stroke clarity.

\section{Production SLURM Batch Workflow and Signal Handling}
\label{sec:slurm_pipeline}

To execute full-page multi-band transcription across entire student cohorts on PARAM Kamrupa, we engineered an automated SLURM batch pipeline with POSIX signal trapping. Because the \texttt{gpu} queue enforces a strict 4-hour walltime limit, long-running batch jobs that do not checkpoint risk sudden termination.

Our production job script (\texttt{scripts/run\_slurm\_slicing.sh}) orchestrates the two-stage execution:
\begin{footnotesize}
\begin{verbatim}
#!/bin/bash
#SBATCH --job-name=btp_hmer_multiband
#SBATCH --partition=gpu
#SBATCH --nodes=1
#SBATCH --ntasks=1
#SBATCH --cpus-per-task=8
#SBATCH --gres=gpu:1
#SBATCH --time=04:00:00
#SBATCH --output=logs/multiband_%j.log

# Signal trap for graceful SLURM checkpointing 120s before timeout
trap 'echo "[SLURM] 120s Warning: Saving checkpoints..."; exit 0' SIGUSR1

# Environment initialization & dynamic glibc preload
module load Anaconda3-2023.03
source activate /scratch/p.dinkesh/hmer_project/venv/
export LD_PRELOAD=/opt/ohpc/pub/apps/Anaconda3-2023.03/lib/libstdc++.so.6

# Stage 1: High-Resolution Multi-Band Slicing & Vision-Language Inference
python batch_slice_transcribe.py \
    --input_dir /scratch/p.dinkesh/hmer_project/data/exam_pages/ \
    --output_dir /scratch/p.dinkesh/hmer_project/results/slices/ \
    --max_dim 1200 --overlap_ratio 0.08

# Stage 2: Pushdown Automaton LaTeX Stitching & Autonomous Tectonic Compilation
python stitch_latex_parts.py \
    --slices_dir /scratch/p.dinkesh/hmer_project/results/slices/ \
    --output_pdf_dir /scratch/p.dinkesh/hmer_project/results/compiled_pdfs/ \
    --compile_engine tectonic
\end{verbatim}
\end{footnotesize}
This two-stage workflow runs completely unattended on compute nodes, transcribing high-resolution horizontal bands, synthesizing balanced \LaTeX\ ASTs via our pushdown automaton, and compiling publication-grade PDFs via Tectonic in under 5 seconds per page.
